\clearpage

% --------------------------------------------------------------------------------------------

\section {Summary}
\label{sec:summary}

In this technical note,
we have presented the results of a 3-flavour $\bar{\nu}_{\mu}$-disappearance analysis,
performed by the {\em VALOR} group,
on the combined Run 5+6 \TotalPot{} dataset, using only data taken in antineutrino beam mode.
This analysis has adapted the work presented in
Refs. \cite{OurRun1DisappAnalysis}, 
\cite{OurRun12DisappAnalysis}, 
\cite{OurRun12DisappAnalysis3Flavour},
\cite{OurRun123DisappAnalysis3Flavour},
\cite{OurRun1234aDisappAnalysis3Flavour} and 
\cite{OurRun1234DisappAnalysis3Flavour}
to this dataset. The Run 5+6 dataset corresponds to an integrated J-PARC neutrino beam
exposure of \TotalPot{} POT.

This analysis predicts \NskNoOscExpectation{} $\pm$ \NskNoOscExpectationError{} (syst)
single $\mu$-like ring events in Super K in the absence of any oscillation,
but only \NskObserved{} were observed. 
The observed deficit has a strong energy dependence: the ratio of the number of observed events to 
the number expected under the no-oscillation hypothesis is
$\sim$\NskRatioLowEnergy\% $<$ 0.5 GeV,
$\sim$\NskRatioMidEnergy\% between 0.5 and 1 GeV and
$\sim$\NskRatioHighEnergy\% $>$ 1 GeV.

One of the motivations for this analysis is to investigate whether the oscillation parameters for $\bar{\nu}_{\mu}$ 
disappearance are different from those for $\nu_{\mu}$ disappearance. In order to allow for this possibility, 
two separate sets of oscillation parameters are used. One set is applied to the $\bar{\nu}_{\mu}$, $\bar{\nu}_{e}$ 
and oscillated $\bar{\nu}_{e}$ interaction modes, while the other set is applied to the $\nu_{\mu}$, $\nu_{e}$ and 
oscillated $\nu_{e}$ modes.

A $\bar{\nu}_{\mu}$-disappearance analysis was performed
in a framework of 3-flavour oscillations including matter effects in constant-density matter.
The observed reconstructed energy spectrum of single $\mu$-like ring events was fitted,
and separate fits were made for the normal and the inverted mass hierarchies. In these fits, $\sin^{2}\overline{\theta}_{23}$ and 
$|\overline{\Delta m}^{2}_{32}|$
%either $|\overline{\Delta m}^{2}_{32}|$ (normal hierarchy) or $|\overline{\Delta m}^{2}_{31}|$ (inverted hierarchy) 
were allowed to float, while all other oscillation
parameters in both sets were fixed to T2K best-fit or PDG values. All \NSyst{} systematic
parameters considered in this analysis were also allowed to float in the fits.

For the normal mass hierarchy, the 3-flavour $\nu_{\mu}$-disappearance fit of the combined Run 5+6 dataset
gives $\sin^{2}\overline{\theta}_{23}$ = \BestFitMixingNormal{} and $|\overline{\Delta m}^{2}_{32}|$ = \BestFitDmsqNormal{}
eV$^{2}$/c$^{4}$ as the best-fit values (p-value = \PValueNormal). The fit for the inverted mass hierarchy yields
$\sin^{2}\overline{\theta}_{23}$ = \BestFitMixingInverted{} and $|\overline{\Delta m}^{2}_{31}|$ = \BestFitDmsqInverted{}
eV$^{2}$/c$^{4}$ as the best-fit values (p-value = \PValueInverted). The p-value for the normal hierarchy antineutrino-mode dataset 
given a null hypothesis of the normal hierarchy best-fit values of the equivalent parameters from the official Run 1-4 
neutrino-mode analysis is \PValueNuModeComp; thus the results of this analysis considered together with the Run 1-4 neutrino-mode results 
are consistent with CPT being conserved.

Confidence regions in the parameter space $\sin^{2}\overline{\theta}_{23}$ - $|\overline{\Delta m}^{2}_{32}|$ 
%(normal hierarchy) and $\sin^{2}\overline{\theta}_{23}$ - $|\overline{\Delta m}^{2}_{31}|$ (inverted hierarchy)
are shown using 
%both 
the constant-$\Delta\chi^2$ method
%and the method of Feldman and Cousins
.
